\documentclass{report}
\usepackage{amsmath,amssymb}
\setlength{\parindent}{0mm}
\setlength{\parskip}{1em}
\begin{document}
\begin{center}
\rule{6in}{1pt} \
{\large Stefan L. Glimberg \\
{\bf Multi-level parallelization via spatial and temporal decomposition on heterogeneous hardware}}

Technical University of Denmark \\ Richard Petersens Plads \\ Building 321/020 \\ 2800 Kgs -Lyngby
\\
{\tt slgl@imm.dtu.dk}\\
Allan P. Engsig-Karup\\
Allan S. Nielsen\end{center}

In this talk we present an efficient solution strategy for fully
nonlinear free surface water waves, based on unified potential flow
theory, with the bottleneck problem of solving a $\sigma$-transformed
Laplace problem in three dimensions at every time integration step. A
geometric multigrid preconditioned defect correction scheme is used to
attain high-order accurate solutions with fast convergence and scalable
work effort. The numerical method is based on matrix-free finite
difference approximations, implemented to run efficiently on many-core
GPUs~\cite{EngsigKarupEtAl2011,GlimbergEtAl2011}. In this talk we present
the extension of previous work, with the addition of both spatial and
temporal decomposition techniques for fast simulation of large scale
phenomenons, such as long distance wave propagation over varying depths
or within large coastal regions. Simulations that have novel value within
maritime engineering because of the tunable properties that follow from
the flexible-order implementation.

The architectural changes in hardware design within the last two decades,
from single to multi- and many-core architectures, require software
developers to identify and properly implement methods that both exploit
concurrency and maintain numerical efficiency. We discuss the challenges
of implementing an effective multigrid solver on modern many-core GPUs
and in particular how multiple devices are used to further improve
performance via distributed heterogeneous computing without compromising
numerical convergence. We present a multi-block approach that decompose
and numerically solve the Laplace problem within each subdomain,
supporting flexible block structures to match the physical domain.
Messages are sent using MPI to repeatedly update artificial boundary
information between adjacent subdomains. The impact on convergence and
performance scalability using the proposed multi-block strategy will be
discussed.

We find that spatial domain decomposition scales well for large problems
sizes, but for problems of limited sizes, the maximum attainable speedup
is reached for a low number of processors, as it leads to an unfavorable
communication to compute ratio. To circumvent this, we exploit the
Parareal algorithm to introduce parallelism via parallel time
integration. Parareal may be perceived as a two level multigrid method in
time, where the numerical solution is first sequentially advanced via
course integration and then updated simultaneously on multiple GPUs in a
predictor-corrector fashion~\cite{LSY02}. A parameter study is performed
to establish proper choices for maximizing speedup and parallel
efficiency. The parareal algorithm is found to be sensitive to a number
of numerical and physical parameters, making practical speedup a matter
of parameter tuning. Results are presented to confirm that it is possible
to attain reasonable speedups, independently of the problem size.

\begin{thebibliography}{50}
\bibitem{EngsigKarupEtAl2011}
Engsig-Karup, A.~P., Madsen, M.~G., Glimberg, S.~L., A massively parallel
gpu-accelerated model for analysis of fully nonlinear free surface waves.
International Journal for Numerical Methods in Fluids 70~(1), 20--36,
2011.
\bibitem{GlimbergEtAl2011}
Glimberg, S.~L., Engsig-Karup, A.~P., Madsen, M.~G., A fast
gpu-accelerated mixed-precision strategy for fully nonlinear water wave
computations. Proceedings of ENUMATH, the 9th European Conference
on Numerical Mathematics and Advanced Applications, 2011.
\bibitem{LSY02}
Baffico, L., Bernard, S., Maday, Y., Turinici, G., Z\'{e}rah, G.,
Parallel in time molecular dynamics simulations. Physical Review E. 66,
2002.
\end{thebibliography}


\end{document}
